% What good annotation looks like, using a sentence from M&L p. 12.
\begin{guidefig}{A sample of good annotation: one key sentence marked, a word defined in the margin, a question, and a connection.}
\begin{tikzpicture}
  \node[draw=black!25,fill=white,text width=8.6cm,align=left,inner sep=8pt,font=\small] (p) at (0,0) {%
    After posing questions, scientists use further observations to make \underline{inferences}.\,\textsuperscript{\textcolor{vocabc}{\textbf{1}}}
    \colorbox{keyc!30}{An inference is a logical interpretation based on}
    \colorbox{keyc!30}{what scientists already know.}\,\textsuperscript{\textcolor{keyc}{\textbf{2}}}
    Inference, combined with a creative imagination, can lead to a hypothesis.\,\textsuperscript{\textcolor{watchc}{\textbf{3}}}
    A hypothesis is a tentative scientific explanation that can be tested by further observation or by experimentation.\,\textsuperscript{\textcolor{linkc}{\textbf{4}}}};
  \node[anchor=north west,text width=5.2cm,align=left,font=\small\sffamily] at ($(p.north east)+(0.4,0)$) {%
    \textcolor{vocabc}{\textbf{1 Define:}} inference = what I \emph{conclude}, not what I see\\[4pt]
    \textcolor{keyc}{\textbf{2 Key sentence:}} highlight only this one\\[4pt]
    \textcolor{watchc}{\textbf{3 Question:}} how is a hypothesis different from an inference?\\[4pt]
    \textcolor{linkc}{\textbf{4 Connect:}} same as the salt-marsh example on p. 11};
\end{tikzpicture}
\end{guidefig}
